%! TEX root = kelengkapan_k3_temu3.tex
% the main TeX file which is intended to compile, :VimtexReload after adjustment
% :h vimtex-tex-root
\documentclass[../main.tex]{subfiles}
% \includeonlyframes{current}
% \setbeameroption{hide notes} % Only slides
% \setbeameroption{show only notes} % Only notes
% \setbeameroption{show notes on second screen=right} % Both

% Identitas Presentasi
\title{Kelengkapan Keselamatan dan Kesehatan Kerja (K3) pada Proyek Konstruksi}
\subtitle{Alat Keselamatan Kerja dan Perangkat Proteksi Proyek}
\author{Mata Kuliah: K3 Konstruksi --- Pertemuan 3}
\institute{Dosen Pengampu}
\date{\today}

\begin{document}

% ---------------------------------------------------
% Slide 1 — Judul
% ---------------------------------------------------
\begin{frame}[c, mytitle=center, mycolor=digiPH_darkblue, dark]
	\titlepage
\end{frame}

% ---------------------------------------------------
% Slide 2 — Capaian dan Tujuan Pembelajaran
% ---------------------------------------------------
\begin{frame}[t, mytitle=standard]
	\frametitle{Capaian dan Tujuan Pembelajaran}
	\begin{itemize}
		\item \textbf{CPMK:}
		      \begin{itemize}
			      \item Mengidentifikasi jenis alat K3
			      \item Menjelaskan fungsi dan dasar regulasi
			      \item Menentukan APD sesuai potensi bahaya
		      \end{itemize}
		      \vspace{0.5cm}
		\item \textbf{Tujuan Operasional:}
		      \begin{itemize}
			      \item Menyebutkan ragam APD \& bagian tubuh yang dilindungi
			      \item Menguraikan dasar hukum kewajiban APD
			      \item Menganalisis penempatan alat K3 (APAR, detektor, alarm, P3K)
		      \end{itemize}
	\end{itemize}
\end{frame}

% ---------------------------------------------------
% Slide 3 — Mengapa K3 Penting di Proyek Konstruksi?
% ---------------------------------------------------
\begin{frame}[c, mytitle=center, mycolor=digiPH_gray]
	\frametitle{Mengapa K3 Penting di Proyek Konstruksi?}
	\begin{itemize}
		\item \textbf{Karakteristik risiko tinggi:} Kerja di ketinggian, paparan alat berat, kebisingan, dan bahan kimia.
		\item \textbf{Dampak kecelakaan kerja:} Menimbulkan korban jiwa, kerugian finansial, hingga keterlambatan penyelesaian proyek.
		\item \textbf{Esensi K3:} K3 harus diterapkan sebagai budaya kerja, bukan sekadar kepatuhan administratif semata.
	\end{itemize}
\end{frame}

% ---------------------------------------------------
% Slide 4 — Dasar Hukum K3 di Indonesia
% ---------------------------------------------------
\begin{frame}[t, mytitle=standard]
	\frametitle{Dasar Hukum K3 di Indonesia}
	\begin{itemize}
		\item \textbf{UU No. 1 Tahun 1970} \\ Tentang Keselamatan Kerja.
		\item \textbf{Permenaker PER.08/MEN/VII/2010} \\ Mengatur kewajiban penyediaan \& penggunaan APD.
		\item \textbf{PP No. 50 Tahun 2012} \\ Penerapan SMK3 di tingkat perusahaan.
		\item \textbf{Permen PUPR No. 10 Tahun 2021} \\ SMKK khusus untuk sektor konstruksi.
		\item \textbf{SNI ISO 45001:2018} \\ Standar internasional sistem manajemen K3 yang diadopsi BSN.
	\end{itemize}
\end{frame}

% ---------------------------------------------------
% Slide 5 — Definisi dan Klasifikasi APD
% ---------------------------------------------------
\begin{frame}[c, mytitle=imageplus]
	\frametitle{Definisi dan Klasifikasi APD}
	\begin{itemize}
		\item \textbf{Definisi APD (Permenaker 08/2010):} \\ Alat yang mempunyai kemampuan untuk mengisolasi sebagian atau seluruh tubuh dari potensi bahaya di tempat kerja.
		\item \textbf{Kelompok APD:} \\ Meliputi pelindung kepala, mata-wajah, telinga, pernapasan, tubuh, tangan, dan kaki.
		\item \textbf{Prinsip Penggunaan:} \\ APD berfungsi sebagai lapisan pertahanan terakhir (last resort) setelah pengendalian teknis dan administratif dilakukan.
	\end{itemize}
\end{frame}

% ---------------------------------------------------
% Slide 6 — Alat Penanda dan Promosi K3
% ---------------------------------------------------
\begin{frame}[t, mytitle=standard, mycolor=digiPH_cream]
	\frametitle{Alat Penanda dan Promosi K3}
	\begin{itemize}
		\item \textbf{Bendera K3}
		      \begin{itemize}
			      \item Dipasang selama aktivitas proyek berlangsung.
			      \item Berfungsi sebagai simbol kesiapsiagaan K3 di lokasi.
		      \end{itemize}
		      \vspace{0.3cm}
		\item \textbf{Sign board K3}
		      \begin{itemize}
			      \item Rambu peringatan, larangan, kewajiban, dan informasi keselamatan.
		      \end{itemize}
		      \vspace{0.3cm}
		\item \textbf{Fungsi Utama:} \\ Membangun kesadaran kolektif para pekerja terhadap berbagai potensi bahaya di sekeliling mereka.
	\end{itemize}
\end{frame}

% ---------------------------------------------------
% Slide 7 — Pelindung Kepala dan Wajah
% ---------------------------------------------------
\begin{frame}[c, mytitle=standard]
	\frametitle{Pelindung Kepala dan Wajah}
	\begin{itemize}
		\item \textbf{Safety Helmet (Helm Keselamatan):} \\ Melindungi kepala dari benturan, kejatuhan benda dari atas, dan pukulan benda keras.
		\item \textbf{Kacamata Safety (Safety Glasses):} \\ Melindungi mata dari paparan debu, partikel berterbangan, percikan bahan kimia, serta cahaya silau.
		\item \textbf{Pelindung Wajah (Face Shield):} \\ Memberikan perlindungan yang lebih menyeluruh untuk wajah dari material berbahaya dan percikan api/kimia.
	\end{itemize}
\end{frame}

% ---------------------------------------------------
% Slide 8 — Pelindung Pernapasan
% ---------------------------------------------------
\begin{frame}[t, mytitle=center, mycolor=digiPH_gray]
	\frametitle{Pelindung Pernapasan}
	\begin{itemize}
		\item \textbf{Masker (Kain/Kertas):}
		      \begin{itemize}
			      \item Melindungi hidung dan mulut dari debu serta bahan kimia ringan.
		      \end{itemize}
		\item \textbf{Respirator:}
		      \begin{itemize}
			      \item Memiliki fungsi serupa dengan masker, namun dirancang khusus untuk lingkungan berisiko tinggi (misalnya terdapat gas beracun).
		      \end{itemize}
		\item \textbf{Kriteria Pemilihan:}
		      \begin{itemize}
			      \item Harus didasarkan pada jenis paparan dan tingkat konsentrasi kontaminan udara di lokasi kerja.
		      \end{itemize}
	\end{itemize}
\end{frame}

% ---------------------------------------------------
% Slide 9 — Pelindung Pendengaran
% ---------------------------------------------------
\begin{frame}[c, mytitle=imageplus]
	\frametitle{Pelindung Pendengaran}
	\begin{itemize}
		\item \textbf{Ear Plug (Sumbat Telinga):}
		      \begin{itemize}
			      \item Dimasukkan ke dalam liang telinga.
			      \item Mampu mereduksi kebisingan sekitar 10–15 dB.
		      \end{itemize}
		      \vspace{0.3cm}
		\item \textbf{Ear Muff (Penutup Telinga):}
		      \begin{itemize}
			      \item Menutup seluruh daun telinga dari luar.
			      \item Mampu mereduksi kebisingan lebih besar, yaitu 20–30 dB.
		      \end{itemize}
		      \vspace{0.3cm}
		\item \textbf{Pemilihan APD:} \\ Wajib disesuaikan dengan intensitas kebisingan aktual di lokasi kerja.
	\end{itemize}
\end{frame}

% ---------------------------------------------------
% Slide 10 — Pelindung Tubuh dan Tangan
% ---------------------------------------------------
\begin{frame}[t, mytitle=standard]
	\frametitle{Pelindung Tubuh dan Tangan}
	\begin{itemize}
		\item \textbf{Wearpack / Coverall:} \\ Pakaian pelindung yang menutup area leher hingga mata kaki, wajib digunakan di lingkungan dengan risiko tinggi.
		\item \textbf{Rompi Safety:} \\ Memiliki material reflektif untuk meningkatkan visibilitas pekerja, sangat penting untuk kerja malam hari atau di area minim pencahayaan.
		\item \textbf{Sarung Tangan (Safety Gloves):} \\ Melindungi tangan dari paparan bahan kimia, sengatan listrik, abrasi, dan benda tajam.
	\end{itemize}
\end{frame}

% ---------------------------------------------------
% Slide 11 — Pelindung Kaki dan Alat Keselamatan Khusus
% ---------------------------------------------------
\begin{frame}[c, mytitle=center]
	\frametitle{Pelindung Kaki dan Alat Keselamatan Khusus}
	\begin{itemize}
		\item \textbf{Sepatu Pelindung (Safety Shoes):} \\ Melindungi kaki dari tertusuk benda tajam, kejatuhan benda berat, tumpahan larutan kimia, hingga sengatan arus listrik.
		\item \textbf{Pelampung (Life Vest):} \\ Mencegah risiko tenggelam pada proyek yang berada di dekat atau di atas wilayah perairan.
		\item \textbf{Sabuk Pengaman (Seat Belt):} \\ Menjadi pelindung wajib bagi pengemudi dan penumpang kendaraan berat di proyek.
		\item \textbf{Body Harness:} \\ Mencegah risiko jatuh fatal saat bekerja di ketinggian, dengan titik tambatan beban di bagian dada/punggung.
	\end{itemize}
\end{frame}

% ---------------------------------------------------
% Slide 12 — Perangkat Proteksi Khusus Proyek
% ---------------------------------------------------
\begin{frame}[t, mytitle=standard, mycolor=digiPH_yellow]
	\frametitle{Perangkat Proteksi Khusus Proyek}
	\begin{itemize}
		\item \textbf{APAR (Alat Pemadam Api Ringan):} Respons dini terhadap risiko kebakaran.
		\item \textbf{Detektor:} Sensor untuk mendeteksi paparan gas beracun (CO), asap, dan suhu panas berlebih dari arus listrik.
		\item \textbf{Alarm:} Memberikan peringatan dini secara audible (terdengar) dari hasil tangkapan detektor.
		\item \textbf{Pelindung Fisik (Physical Guard):} Pagar atau sekat untuk mencegah kontak langsung antara pekerja dengan sumber bahaya (misal: mesin berputar).
		\item \textbf{Kotak P3K:} Menunjang pertolongan pertama secara cepat saat terjadi insiden kecelakaan kerja.
	\end{itemize}
\end{frame}

% ---------------------------------------------------
% Slide 13 — Studi Kasus: Kecelakaan Akibat Kelalaian
% ---------------------------------------------------
\begin{frame}[c, mytitle=standard, mycolor=digiPH_ocean, dark]
	\frametitle{Studi Kasus: Kecelakaan Akibat Kelalaian APD}
	\begin{itemize}
		\item \textbf{Paparan Kasus Nyata:} \\ Analisis kejadian kecelakaan kerja di sektor konstruksi (seperti insiden jatuh dari ketinggian atau kebakaran area kerja).
		\item \textbf{Diskusi Kelompok:}
		      \begin{itemize}
			      \item Mengidentifikasi jenis APD yang seharusnya wajib digunakan.
			      \item Menganalisis akar penyebab utama (root cause) kelalaian.
		      \end{itemize}
		\item \textbf{Aspek Hukum:} \\ Kaitan kelalaian dalam kasus ini dengan ancaman sanksi bagi pihak pengusaha sesuai dengan Permenaker 08/2010.
	\end{itemize}
\end{frame}

% ---------------------------------------------------
% Slide 14 — Latihan: Menentukan Kombinasi APD
% ---------------------------------------------------
\begin{frame}[t, mytitle=imageplus]
	\frametitle{Latihan: Menentukan Kombinasi APD}
	\textbf{Tugas Mahasiswa:} Susun daftar APD beserta justifikasi untuk skenario berikut:
	\vspace{0.3cm}
	\begin{itemize}
		\item \textbf{Skenario 1 - Kerja di Ketinggian:}
		      \begin{itemize}
			      \item Kebutuhan dasar: Body harness, helm, sepatu anti-slip.
		      \end{itemize}
		\item \textbf{Skenario 2 - Area Bising Tinggi:}
		      \begin{itemize}
			      \item Kebutuhan dasar: Ear muff/ear plug, sistem komunikasi alternatif.
		      \end{itemize}
		\item \textbf{Skenario 3 - Ruang Berisiko Kebakaran:}
		      \begin{itemize}
			      \item Kebutuhan dasar: APAR stand-by, detektor, alarm, wearpack tahan api.
		      \end{itemize}
	\end{itemize}
\end{frame}

% ---------------------------------------------------
% Slide 15 — Penataan Alat K3 dan Penutup
% ---------------------------------------------------
\begin{frame}[c, mytitle=center]
	\frametitle{Penataan Alat K3 dan Penutup}
	\begin{itemize}
		\item \textbf{Prinsip Penempatan Alat K3 (APAR, P3K, dll):}
		      \begin{itemize}
			      \item Harus mudah dijangkau dan terlihat jelas oleh siapa saja.
			      \item Ditempatkan dekat area dengan potensi risiko insiden tinggi.
			      \item \textbf{Tidak boleh} menghalangi jalur evakuasi.
		      \end{itemize}
		      \vspace{0.4cm}
		\item \textbf{Refleksi Akhir:} \\ Membangun budaya keselamatan kerja secara berkelanjutan adalah tanggung jawab bersama di lingkungan proyek.
		      \vspace{0.4cm}
		\item \textbf{Sesi Penutup:} \\ Tanya jawab \& rangkuman keseluruhan capaian pembelajaran hari ini.
	\end{itemize}
\end{frame}

\end{document}
